\chapter*{Avant-propos}
\addcontentsline{toc}{chapter}{Avant-propos}
\markboth{Avant-propos}{}

\begin{quote}
\itshape « Un article de recherche ne devrait effrayer personne. »
\end{quote}

Ce livre raconte un seul article scientifique, du premier mot jusqu'à la
dernière ligne de code : \emph{V-JEPA~2.1: Unlocking Dense Features in Video
Self-Supervised Learning}~\parencite{murlabadia2026vjepa21}, écrit par
l'équipe FAIR de Meta et l'Université de Saragosse. Cet article présente un
programme qui apprend à comprendre des vidéos tout seul, sans qu'un humain
lui explique jamais ce qu'il regarde. Il fait suite à un premier modèle,
V-JEPA~2~\parencite{assran2025vjepa2}, dont nous parlerons beaucoup : c'est
en corrigeant ses défauts que V-JEPA~2.1 est né.

Un article de recherche est écrit par des spécialistes, pour des
spécialistes. Il va vite. Il emploie des mots que personne ne vous a jamais
définis. Il suppose que vous connaissez déjà dix autres articles. Résultat :
beaucoup de gens curieux, capables, motivés, s'arrêtent à la première page.

\textbf{Ici, on ne suppose rien du tout.} Chaque idée est d'abord racontée
avec des mots de tous les jours, puis dessinée, puis seulement écrite en
langage mathématique, puis enfin mise en pratique avec du code que vous
pouvez faire tourner vous-même.

\begin{enclair}
Notre promesse : si vous savez compter et si vous savez lire, vous pouvez
comprendre comment fonctionne V-JEPA~2.1. Le reste, nous le construisons
ensemble, dans l'ordre, sans jamais sauter une marche.
\end{enclair}

\section*{Pour qui est ce livre}

Pour toute personne qui veut trois choses à la fois :

\begin{enumerate}
  \item \textbf{comprendre} ce que fait V-JEPA~2.1, et pourquoi il fait
        mieux que le modèle qui existait avant lui ;
  \item \textbf{écrire elle-même} le modèle et le programme qui l'entraîne ;
  \item \textbf{mettre le modèle au travail} pour de vrai, tous les jours,
        au service d'une application.
\end{enumerate}

Aucun diplôme n'est demandé. Vous n'avez pas besoin de savoir ce qu'est un
réseau de neurones, ni un transformeur, ni même une image numérique : tout
cela est expliqué à partir de rien, dans les quatre premiers chapitres.

\section*{Comment ce livre est fabriqué}

Huit règles ont guidé chaque page, et elles ne sont pas des intentions : ce
sont des engagements que vous pouvez vérifier.

\begin{enumerate}
  \item Aucun mot technique n'apparaît sans avoir été défini juste avant.
  \item On commence toujours par les prérequis, jamais par la nouveauté.
  \item Un dessin pour chaque idée qui bouge, qui circule ou qui se compare.
  \item Un exemple après chaque notion, avec des chiffres quand c'est
        possible.
  \item Des mots simples, ceux de tous les jours.
  \item Le texte vous parle et vous accompagne.
  \item Une phrase de transition à la fin de chaque section, qui annonce la
        suite.
  \item Des phrases entières, comme dans un vrai livre.
\end{enumerate}

Le texte utilise quelques encadrés qui reviennent souvent, faciles à
reconnaître à leur pictogramme :

\begin{itemize}
  \item \textbf{\faIcon{comments}~En clair} — l'idée expliquée avec des
        mots de tous les jours.
  \item \textbf{\faIcon{brain}~L'intuition} — l'image mentale à garder en
        tête.
  \item \textbf{\faIcon{robot}~À quoi ça sert en IA ?} — où cette notion
        sert vraiment dans V-JEPA~2.1.
  \item \textbf{\faIcon{chart-line}~Sur des données réelles} — la même idée,
        appliquée aux chiffres réels de l'article.
  \item \textbf{\faIcon{seedling}~Ce qu'il faut savoir avant} — les seuls
        prérequis de la section.
  \item \textbf{\faIcon{exclamation-triangle}~Piège à éviter} — l'erreur
        classique que tout le monde fait.
  \item \textbf{\faIcon{key}~À retenir} — l'essentiel, si vous ne deviez
        garder qu'une chose.
  \item \textbf{\faIcon{rocket}~Pour aller plus loin} — un pont vers un
        sujet plus avancé, que vous pouvez sauter en première lecture.
\end{itemize}

Le code, lui, se présente de deux façons, et il ne faut pas les confondre.
Ce que \emph{vous} écrivez apparaît sur fond clair, avec ses numéros de
ligne. Ce que la \emph{machine} répond apparaît dans une petite fenêtre
noire, comme celle d'un vrai terminal, avec ses trois pastilles de couleur
en haut à gauche. Vous saurez ainsi, d'un coup d'œil et sans lire le titre,
si vous avez sous les yeux une instruction à taper ou une réponse à
comparer avec la vôtre.

Chaque chapitre se termine par des exercices, tous corrigés. Ils ne servent
pas à vous noter : ils servent à vous montrer, à vous-même, que vous avez
compris.

\section*{Un mot d'honnêteté}

L'article explique comment \emph{entraîner} le modèle et comment
l'\emph{évaluer}. Il ne parle presque pas de sa \emph{mise en service} : pas
de serveur, pas de temps de réponse, pas de surveillance. La dernière partie
de ce livre sera donc de deux couleurs, et nous vous le dirons franchement à
chaque fois :

\begin{itemize}
  \item \textbf{Ce que l'article dit} — la taille des modèles, la résolution
        des images, la vitesse de calcul. Ce sont des résultats mesurés.
  \item \textbf{Ce que nous ajoutons} — le serveur, la mesure du temps de
        réponse, la surveillance. Ce sont des conseils de métier, utiles et
        éprouvés, mais qui ne viennent pas de l'article.
\end{itemize}

Nous dirons aussi, aussi clairement que ses réussites, ce que ce modèle ne
sait pas faire, et les tâches sur lesquelles d'autres modèles restent
meilleurs que lui. Un livre qui ne raconte que les victoires ne vous aide
pas à choisir.

\vspace{1em}
\hfill\textit{Bonne lecture.}
\clearpage
